\documentclass[12pt]{article}
\usepackage[utf8]{inputenc}
\usepackage[T1]{fontenc}
\usepackage{hyperref}
\title{Yongbei Zhuang Language Resources for K-12 Curriculum Development}
\author{Compiled by Claude Sonnet 5.5 (claude-sonnet-5-5) for Andr\'as Kornai}
\date{}
\begin{document}
\maketitle

\section{Introduction}
Yongbei Zhuang (ISO 639-3 zyb, Glottocode yong1276) is a Northern Tai variety of the Zhuang macrolanguage, spoken mainly in Guangxi, China. Search results describe its dialects as including Yongning North, Binyang, Hengxian and Pingguo. Wikipedia gives about 2.0 million speakers. Joshua Project estimates 2,155,000. Sources disagree on whether it is grouped under Northern Zhuang or Yongnan-Wuming Zhuang. The search found no school curricula, primers or textbooks specific to this variety, so little K-12 material is documented here.

\section{Primary Educational Resources}
No primers, textbooks or school programmes specific to Yongbei Zhuang turned up in the search results. This section is therefore empty.

\section{Linguistic and Reference Materials}
\begin{itemize}
\item Wikipedia overview of Yongbei Zhuang: \url{https://en.wikipedia.org/wiki/Yongbei_Zhuang}
\item Glottolog language entry (ISO zyb): \url{https://glottolog.org/resource/languoid/iso/zyb}
\item Max Planck EVA numerals page, Zhuang Yongbei North: \url{https://lingweb.eva.mpg.de/channumerals/Zhuang-Yongbei-North.htm}
\item Wikitongues language page: \url{https://wikitongues.org/languages/zyb/}
\end{itemize}

\section{Religious and Community Resources}
\begin{itemize}
\item Joshua Project profile of Yongbei Zhuang in China: \url{https://joshuaproject.net/people_groups/19801}
\item Global Recordings Network page, listed as ``Zhuang, Northern'': \url{https://globalrecordings.net/en/language/zyb}
\end{itemize}

\section{Digital and Audio Resources}
\begin{itemize}
\item Global Recordings Network audio materials: \url{https://globalrecordings.net/en/language/zyb}
\item Detect Language entry: \url{https://www.detectlanguage.com/languages/zyb}
\end{itemize}

\section{Curriculum Development Context}
The language is vigorous. Wikitongues lists it as EGIDS 6a, and Glottolog treats it as not endangered. Its areal group includes the Wuming dialect, which is the basis of Standard Zhuang. Standard Zhuang materials may therefore be a practical bridge, but this report did not locate or verify such materials. Curriculum work will probably need new, locally developed material.

\end{document}